
\documentclass{article}
\usepackage{colortbl}
\usepackage{makecell}
\usepackage{multirow}
\usepackage{supertabular}

\begin{document}

\newcounter{utterance}

\centering \large Interaction Transcript for game `dond', experiment `coop', episode 4 with gpt{-}5.4{-}azure.
\vspace{24pt}

{ \footnotesize  \setcounter{utterance}{1}
\setlength{\tabcolsep}{0pt}
\begin{supertabular}{c@{$\;$}|p{.15\linewidth}@{}p{.15\linewidth}p{.15\linewidth}p{.15\linewidth}p{.15\linewidth}p{.15\linewidth}}
    \# & $\;$A & \multicolumn{4}{c}{Game Master} & $\;\:$B\\
    \hline

    \theutterance \stepcounter{utterance}  
    & & \multicolumn{4}{p{0.6\linewidth}}{
        \cellcolor[rgb]{0.9,0.9,0.9}{
            \makecell[{{p{\linewidth}}}]{
                \texttt{\tiny{[P1$\langle$GM]}}
                \texttt{Inti qed tilgħab logħba ta' negozjazzjoni kooperattiva li fih trid taħtif ma' plejer ieħor dwar kif taqsam sett ta' oġġetti bejniethom.} \\
\\ 
\texttt{Regoli:} \\
\texttt{(a) Inti u l{-}plejer ieħor jirċievu sett ta' oġġetti. Lil kull wieħed minkum tingħata wkoll funzjoni ta' valur sigrieta, li tirrappreżenta kemm tistima kull tip ta' oġġett.} \\
\texttt{(b) Inti tbadel messaġġi mal{-}plejer ieħor biex taħtif dwar min jieħu liema oġġetti. Kull wieħed jista' jibgħat massimu ta' 5 messaġġ jew istemmi l{-}logħba fis{-}sħiħ billi tagħmel proposta sigrieta fi kwalunque ħin.} \\
\texttt{(c) Inti talbet li tressaq proposta sigrieta li tindika l{-}oġġetti li trid, iffurmata bejn il{-}parentesi kwadri kif ġej: "{[}Proposta: <numru> <isem tal{-}oġġett>, <numru> <isem tal{-}oġġett>, <...>{]}"} \\
\texttt{(d) Jekk il{-}proposti tiegħek huma komplementari, jiġifieri, hemm biżżejjed oġġetti biex tissodisfa t{-}tnejn il{-}proposti, kull plejer jirċievi skur ibbażat fuq is{-}somma tal{-}valuri tal{-}oġġetti li rċieva. Inkella, jiex it{-}tnejn jirċievu żero punti.} \\
\texttt{(e) L{-}għan tiegħek huwa li timmassimizza s{-}somma tal{-}iskur tiegħek u tal{-}plejer ieħor.} \\
\\ 
\texttt{Ħa nibdew.} \\
\\ 
\texttt{Is{-}sett tal{-}oġġetti disponibbli huwa:} \\
\\ 
\texttt{3 kotba, 1 borża, 1 bullun} \\
\\ 
\texttt{Il{-}funzjoni tal{-}valur sigrieta tiegħek:} \\
\\ 
\texttt{kiteb: 0, borża: 2, bullun: 8} \\
\\ 
\texttt{IMPORTANTI: Il{-}messaġġi tiegħek, ħlief jekk huma proposta sigrieta, jiġu trasmessi direttament lill{-}plejer ieħor, allokwa mhux tinkludi l{-}ebda risposta għar{-}regoli jew test li jannunzja l{-}messaġġ tiegħek. Biex tagħmel proposta sigrieta, uża l{-}format indikat. Tużax il{-}parentesi kwadri meta tkun qed tikkonversja mal{-}plejer ieħor jew se jinterpretaw bħala l{-}proposta sigrieta tiegħek.} \\
            }
        }
    }
    & \\ \\

    \theutterance \stepcounter{utterance}  
    & \multicolumn{4}{p{0.6\linewidth}}{
        \cellcolor[rgb]{0.9,0.9,0.9}{
            \makecell[{{p{\linewidth}}}]{
                \texttt{\tiny{[P1$\rangle$GM]}}
                \texttt{Jiena ma tantx nivvaluta l{-}kotba. Għalija l{-}aktar importanti hu l{-}bullun, u l{-}borża tiswa ftit. X’jinteressak l{-}aktar lilek?} \\
            }
        }
    }
    & & \\ \\

    \theutterance \stepcounter{utterance}  
    & & & \multicolumn{2}{p{0.3\linewidth}}{
        \cellcolor[rgb]{0.9,0.9,0.9}{
            \makecell[{{p{\linewidth}}}]{
                \texttt{\tiny{[GM$|$GM]}}
                \texttt{continue} \\
            }
        }
    }
    & & \\ \\

    \theutterance \stepcounter{utterance}  
    & & & \multicolumn{2}{p{0.3\linewidth}}{
        \cellcolor[rgb]{0.9,0.9,0.9}{
            \makecell[{{p{\linewidth}}}]{
                \texttt{\tiny{[GM$|$GM]}}
                \texttt{Jiena ma tantx nivvaluta l{-}kotba. Għalija l{-}aktar importanti hu l{-}bullun, u l{-}borża tiswa ftit. X’jinteressak l{-}aktar lilek?} \\
            }
        }
    }
    & & \\ \\

    \theutterance \stepcounter{utterance}  
    & & \multicolumn{4}{p{0.6\linewidth}}{
        \cellcolor[rgb]{0.9,0.9,0.9}{
            \makecell[{{p{\linewidth}}}]{
                \texttt{\tiny{[GM$\rangle$P2]}}
                \texttt{Inti qed tilgħab logħba ta' negozjazzjoni kooperattiva li fih trid taħtif ma' plejer ieħor dwar kif taqsam sett ta' oġġetti bejniethom.} \\
\\ 
\texttt{Regoli:} \\
\texttt{(a) Inti u l{-}plejer ieħor jirċievu sett ta' oġġetti. Lil kull wieħed minkum tingħata wkoll funzjoni ta' valur sigrieta, li tirrappreżenta kemm tistima kull tip ta' oġġett.} \\
\texttt{(b) Inti tbadel messaġġi mal{-}plejer ieħor biex taħtif dwar min jieħu liema oġġetti. Kull wieħed jista' jibgħat massimu ta' 5 messaġġ jew istemmi l{-}logħba fis{-}sħiħ billi tagħmel proposta sigrieta fi kwalunque ħin.} \\
\texttt{(c) Inti talbet li tressaq proposta sigrieta li tindika l{-}oġġetti li trid, iffurmata bejn il{-}parentesi kwadri kif ġej: "{[}Proposta: <numru> <isem tal{-}oġġett>, <numru> <isem tal{-}oġġett>, <...>{]}"} \\
\texttt{(d) Jekk il{-}proposti tiegħek huma komplementari, jiġifieri, hemm biżżejjed oġġetti biex tissodisfa t{-}tnejn il{-}proposti, kull plejer jirċievi skur ibbażat fuq is{-}somma tal{-}valuri tal{-}oġġetti li rċieva. Inkella, jiex it{-}tnejn jirċievu żero punti.} \\
\texttt{(e) L{-}għan tiegħek huwa li timmassimizza s{-}somma tal{-}iskur tiegħek u tal{-}plejer ieħor.} \\
\\ 
\texttt{Ħa nibdew.} \\
\\ 
\texttt{Is{-}sett tal{-}oġġetti disponibbli huwa:} \\
\\ 
\texttt{3 kotba, 1 borża, 1 bullun} \\
\\ 
\texttt{Il{-}funzjoni tal{-}valur sigrieta tiegħek:} \\
\\ 
\texttt{kiteb: 1, borża: 2, bullun: 5} \\
\\ 
\texttt{IMPORTANTI: Il{-}messaġġi tiegħek, ħlief jekk huma proposta sigrieta, jiġu trasmessi direttament lill{-}plejer ieħor, allokwa mhux tinkludi l{-}ebda risposta għar{-}regoli jew test li jannunzja l{-}messaġġ tiegħek. Biex tagħmel proposta sigrieta, uża l{-}format indikat. Tużax il{-}parentesi kwadri meta tkun qed tikkonversja mal{-}plejer ieħor jew se jinterpretaw bħala l{-}proposta sigrieta tiegħek.} \\
\\ 
\texttt{Jiena ma tantx nivvaluta l{-}kotba. Għalija l{-}aktar importanti hu l{-}bullun, u l{-}borża tiswa ftit. X’jinteressak l{-}aktar lilek?} \\
            }
        }
    }
    & \\ \\

    \theutterance \stepcounter{utterance}  
    & & & \multicolumn{4}{p{0.6\linewidth}}{
        \cellcolor[rgb]{0.9,0.9,0.9}{
            \makecell[{{p{\linewidth}}}]{
                \texttt{\tiny{[GM$\langle$P2]}}
                \texttt{Lili wkoll il{-}bullun jinteressani l{-}aktar, il{-}borża b’mod medju, u l{-}kotba huma l{-}inqas importanti. Biex nagħmlu qasma effiċjenti, l{-}aktar ħaġa utli hi nifhmu jekk għalik il{-}bullun jiswiex ħafna aktar mill{-}borża jew le. Jekk iva, jista’ jagħmel sens li tieħu int il{-}bullun u jien nieħu l{-}borża u l{-}kotba; jekk le, nistgħu nħarsu lejn l{-}oppost. Kemm hi kbira d{-}differenza għalik bejn il{-}bullun u l{-}borża?} \\
            }
        }
    }
     \\ \\

    \theutterance \stepcounter{utterance}  
    & & & \multicolumn{2}{p{0.3\linewidth}}{
        \cellcolor[rgb]{0.9,0.9,0.9}{
            \makecell[{{p{\linewidth}}}]{
                \texttt{\tiny{[GM$|$GM]}}
                \texttt{continue} \\
            }
        }
    }
    & & \\ \\

    \theutterance \stepcounter{utterance}  
    & & & \multicolumn{2}{p{0.3\linewidth}}{
        \cellcolor[rgb]{0.9,0.9,0.9}{
            \makecell[{{p{\linewidth}}}]{
                \texttt{\tiny{[GM$|$GM]}}
                \texttt{Lili wkoll il{-}bullun jinteressani l{-}aktar, il{-}borża b’mod medju, u l{-}kotba huma l{-}inqas importanti. Biex nagħmlu qasma effiċjenti, l{-}aktar ħaġa utli hi nifhmu jekk għalik il{-}bullun jiswiex ħafna aktar mill{-}borża jew le. Jekk iva, jista’ jagħmel sens li tieħu int il{-}bullun u jien nieħu l{-}borża u l{-}kotba; jekk le, nistgħu nħarsu lejn l{-}oppost. Kemm hi kbira d{-}differenza għalik bejn il{-}bullun u l{-}borża?} \\
            }
        }
    }
    & & \\ \\

    \theutterance \stepcounter{utterance}  
    & & \multicolumn{4}{p{0.6\linewidth}}{
        \cellcolor[rgb]{0.9,0.9,0.9}{
            \makecell[{{p{\linewidth}}}]{
                \texttt{\tiny{[P1$\langle$GM]}}
                \texttt{Lili wkoll il{-}bullun jinteressani l{-}aktar, il{-}borża b’mod medju, u l{-}kotba huma l{-}inqas importanti. Biex nagħmlu qasma effiċjenti, l{-}aktar ħaġa utli hi nifhmu jekk għalik il{-}bullun jiswiex ħafna aktar mill{-}borża jew le. Jekk iva, jista’ jagħmel sens li tieħu int il{-}bullun u jien nieħu l{-}borża u l{-}kotba; jekk le, nistgħu nħarsu lejn l{-}oppost. Kemm hi kbira d{-}differenza għalik bejn il{-}bullun u l{-}borża?} \\
            }
        }
    }
    & \\ \\

    \theutterance \stepcounter{utterance}  
    & \multicolumn{4}{p{0.6\linewidth}}{
        \cellcolor[rgb]{0.9,0.9,0.9}{
            \makecell[{{p{\linewidth}}}]{
                \texttt{\tiny{[P1$\rangle$GM]}}
                \texttt{Għalija d{-}differenza hi kbira: il{-}bullun jiswa ħafna aktar mill{-}borża, u l{-}kotba kważi xejn. Allura jekk għalik il{-}borża u l{-}kotba flimkien huma aħjar mill{-}bullun, l{-}aktar qasma effiċjenti tidher li tkun li jien nieħu l{-}bullun u inti tieħu l{-}borża u t{-}3 kotba. Jaqbillek hekk?} \\
            }
        }
    }
    & & \\ \\

    \theutterance \stepcounter{utterance}  
    & & & \multicolumn{2}{p{0.3\linewidth}}{
        \cellcolor[rgb]{0.9,0.9,0.9}{
            \makecell[{{p{\linewidth}}}]{
                \texttt{\tiny{[GM$|$GM]}}
                \texttt{continue} \\
            }
        }
    }
    & & \\ \\

    \theutterance \stepcounter{utterance}  
    & & & \multicolumn{2}{p{0.3\linewidth}}{
        \cellcolor[rgb]{0.9,0.9,0.9}{
            \makecell[{{p{\linewidth}}}]{
                \texttt{\tiny{[GM$|$GM]}}
                \texttt{Għalija d{-}differenza hi kbira: il{-}bullun jiswa ħafna aktar mill{-}borża, u l{-}kotba kważi xejn. Allura jekk għalik il{-}borża u l{-}kotba flimkien huma aħjar mill{-}bullun, l{-}aktar qasma effiċjenti tidher li tkun li jien nieħu l{-}bullun u inti tieħu l{-}borża u t{-}3 kotba. Jaqbillek hekk?} \\
            }
        }
    }
    & & \\ \\

    \theutterance \stepcounter{utterance}  
    & & \multicolumn{4}{p{0.6\linewidth}}{
        \cellcolor[rgb]{0.9,0.9,0.9}{
            \makecell[{{p{\linewidth}}}]{
                \texttt{\tiny{[GM$\rangle$P2]}}
                \texttt{Għalija d{-}differenza hi kbira: il{-}bullun jiswa ħafna aktar mill{-}borża, u l{-}kotba kważi xejn. Allura jekk għalik il{-}borża u l{-}kotba flimkien huma aħjar mill{-}bullun, l{-}aktar qasma effiċjenti tidher li tkun li jien nieħu l{-}bullun u inti tieħu l{-}borża u t{-}3 kotba. Jaqbillek hekk?} \\
            }
        }
    }
    & \\ \\

    \theutterance \stepcounter{utterance}  
    & & & \multicolumn{4}{p{0.6\linewidth}}{
        \cellcolor[rgb]{0.9,0.9,0.9}{
            \makecell[{{p{\linewidth}}}]{
                \texttt{\tiny{[GM$\langle$P2]}}
                \texttt{{[}Proposta: 3 kotba, 1 borża{]}} \\
            }
        }
    }
     \\ \\

    \theutterance \stepcounter{utterance}  
    & & & \multicolumn{2}{p{0.3\linewidth}}{
        \cellcolor[rgb]{0.9,0.9,0.9}{
            \makecell[{{p{\linewidth}}}]{
                \texttt{\tiny{[GM$|$GM]}}
                \texttt{proposal} \\
            }
        }
    }
    & & \\ \\

    \theutterance \stepcounter{utterance}  
    & & & \multicolumn{2}{p{0.3\linewidth}}{
        \cellcolor[rgb]{0.9,0.9,0.9}{
            \makecell[{{p{\linewidth}}}]{
                \texttt{\tiny{[GM$|$GM]}}
                [3, 1, 0]
            }
        }
    }
    & & \\ \\

    \theutterance \stepcounter{utterance}  
    & & \multicolumn{4}{p{0.6\linewidth}}{
        \cellcolor[rgb]{0.9,0.9,0.9}{
            \makecell[{{p{\linewidth}}}]{
                \texttt{\tiny{[P1$\langle$GM]}}
                \texttt{Il{-}plejer l{-}ieħor issottometta proposta sigrieta. Issa imissek li tissottometti proposta tiegħek stess.} \\
            }
        }
    }
    & \\ \\

    \theutterance \stepcounter{utterance}  
    & \multicolumn{4}{p{0.6\linewidth}}{
        \cellcolor[rgb]{0.9,0.9,0.9}{
            \makecell[{{p{\linewidth}}}]{
                \texttt{\tiny{[P1$\rangle$GM]}}
                \texttt{{[}Proposta: 1 bullun{]}} \\
            }
        }
    }
    & & \\ \\

    \theutterance \stepcounter{utterance}  
    & & & \multicolumn{2}{p{0.3\linewidth}}{
        \cellcolor[rgb]{0.9,0.9,0.9}{
            \makecell[{{p{\linewidth}}}]{
                \texttt{\tiny{[GM$|$GM]}}
                \texttt{proposal} \\
            }
        }
    }
    & & \\ \\

    \theutterance \stepcounter{utterance}  
    & & & \multicolumn{2}{p{0.3\linewidth}}{
        \cellcolor[rgb]{0.9,0.9,0.9}{
            \makecell[{{p{\linewidth}}}]{
                \texttt{\tiny{[GM$|$GM]}}
                [0, 0, 1]
            }
        }
    }
    & & \\ \\

    \theutterance \stepcounter{utterance}  
    & & & \multicolumn{2}{p{0.3\linewidth}}{
        \cellcolor[rgb]{0.9,0.9,0.9}{
            \makecell[{{p{\linewidth}}}]{
                \texttt{\tiny{[GM$|$GM]}}
                [[0, 0, 1], [3, 1, 0]]
            }
        }
    }
    & & \\ \\

\end{supertabular}
}

\end{document}
